\documentclass[11pt]{amsart}
\usepackage[margin=1.1in]{geometry}
\usepackage{amssymb}
\usepackage[T1]{fontenc}
\usepackage{lmodern}
\usepackage{listings}
\usepackage{xcolor}
\usepackage[colorlinks=true,linkcolor=blue,citecolor=blue,urlcolor=blue]{hyperref}
\usepackage[capitalize,nameinlink]{cleveref}

\lstdefinelanguage{Lean}{
  morekeywords={def,theorem,lemma,structure,where,fun,Prop,Type,by,import,namespace,end,noncomputable,instance},
  sensitive=true, morecomment=[l]{--}, morecomment=[s]{/-}{-/}, morestring=[b]"}
\lstset{language=Lean, basicstyle=\ttfamily\small, keywordstyle=\bfseries\color{blue!60!black},
  commentstyle=\itshape\color{gray}, columns=fullflexible, breaklines=true, frame=single,
  extendedchars=true, inputencoding=utf8}

\theoremstyle{plain}
\newtheorem{theorem}{Theorem}[section]
\newtheorem{lemma}[theorem]{Lemma}
\newtheorem{proposition}[theorem]{Proposition}
\newtheorem{corollary}[theorem]{Corollary}
\theoremstyle{definition}
\newtheorem{definition}[theorem]{Definition}
\newtheorem{hypothesis}[theorem]{Hypothesis}
\theoremstyle{remark}
\newtheorem{remark}[theorem]{Remark}

\newcommand{\DHL}{\mathrm{DHL}}

\title[DHL\texorpdfstring{$[3,2]$}{[3,2]} yields \texorpdfstring{$H_1\le 6$}{H1<=6}, not \texorpdfstring{$H_1=2$}{H1=2}]{Resolving a category error: \texorpdfstring{$\mathrm{DHL}[3,2]$}{DHL[3,2]} yields \texorpdfstring{$H_1 \le 6$}{H1 <= 6}, not \texorpdfstring{$H_1=2$}{H1 = 2}}
\author{Benjamin Frohman}
\email{frohmanbenjamin@gmail.com}
\subjclass[2020]{Primary 11N05; Secondary 11N35, 03B35}
\keywords{prime gaps, bounded gaps between primes, admissible tuples, Dickson--Hardy--Littlewood conjecture, parity barrier, twin prime conjecture, formal verification, Lean 4, Mathlib}
\date{November 2026}
\emergencystretch=3em


\begin{document}
\begin{abstract}
Let $H_1=\liminf_{n\to\infty}(p_{n+1}-p_n)$. We record the precise logical relationship between the hypothesis $\DHL[3,2]$ and bounds on $H_1$. We prove that $\DHL[3,2]$ implies $H_1\le 6$; that $H_1\le 3$ is equivalent to $H_1=2$ (the twin prime conjecture); that the minimal diameter of an admissible $3$-tuple is exactly $6$; and, via an explicit model, that $\DHL[2,2]$ cannot be derived from the combinatorial form of $\DHL[3,2]$ alone. The parity barrier for square-of-divisor-sum sieves enters only as an explicitly stated hypothesis. We do not prove $\DHL[3,2]$ or the twin prime conjecture. All results are formalized in Lean~4 with Mathlib, without \texttt{sorry} and using only Lean's standard axioms.
\end{abstract}
\maketitle

\section{Introduction}\label{sec:intro}

Let $p_1=2<p_2=3<\dots$ denote the primes and $H_1=\liminf_{n\to\infty}(p_{n+1}-p_n)$. For an admissible $k$-tuple $\mathcal H$ (\cref{def:admissible}), the statement $\DHL[k,j]$ asserts that infinitely many translates $n+\mathcal H$ contain at least $j$ primes (\cref{def:dhl}).

Following Goldston, Pintz and Y{\i}ld{\i}r{\i}m~\cite{GPY}, Maynard~\cite{Maynard15} and the Polymath8b project~\cite{Polymath}, the unconditional bound $H_1\le 246$ is known, and $\DHL[3,2]$ is known conditionally on the generalized Elliott--Halberstam conjecture (GEH), yielding $H_1\le 6$ under GEH~\cite{Polymath}. A claimed improvement to $H_1\le 240$, attributed to a 2026 preprint of Stadlmann, is \textbf{[citation needed: unverified]} and is not used here.

The purpose of this note is to separate two statements that are sometimes conflated: $\DHL[3,2]\Rightarrow H_1\le 6$ (\cref{thm:main}), which is elementary, and $H_1=2$, the twin prime conjecture. Since every prime gap after $3-2$ is even (\cref{thm:gaps}), $H_1\le 3$ is equivalent to $H_1=2$, so reading $\DHL[3,2]$ as giving $H_1\le 3$ is a category error.

By the widely accepted parity barrier (going back to Selberg), standard square-of-divisor-sum sieves are not expected to go below $6$, even under GEH; in this note that barrier is used only as an explicit hypothesis (\cref{hyp:parity}), not proved.

\section{Preliminaries}\label{sec:prelim}

\begin{definition}\label{def:admissible}
A finite set $\mathcal H\subset\mathbb N$ is \emph{admissible} if for every prime $q$ there is a residue class modulo $q$ containing no element of $\mathcal H$.
\end{definition}

\begin{definition}\label{def:dhl}
For a predicate $P$ on $\mathbb N$ and $k\ge j\ge1$, let $\DHL_P[k,j]$ be the statement: for every admissible $\mathcal H$ with $|\mathcal H|=k$ there are infinitely many $n$ such that $\#\{h\in\mathcal H: P(n+h)\}\ge j$. Write $\DHL[k,j]$ for $\DHL_P[k,j]$ with $P$ = primality.
\end{definition}

\begin{definition}\label{def:H1le}
For $c\in\mathbb N$ write $H_1\le c$ for the statement that $p_{n+1}-p_n\le c$ for infinitely many $n$. For an integer-valued sequence this is equivalent to $\liminf_{n}(p_{n+1}-p_n)\le c$.
\end{definition}

\begin{lemma}\label{lem:adm}
The sets $\{0,2\}$, $\{0,2,6\}$ and $\{0,4,6\}$ are admissible.
\end{lemma}
\begin{proof}
For $q\ge7$ the elements are distinct residues below $q$ and the class $1$ is missed. For $q=2$ the class $1$ is missed. For $q=3$: $\{0,4,6\}\equiv\{0,1,0\}$ misses $2$; $\{0,2,6\}\equiv\{0,2,0\}$ and $\{0,2\}$ miss $1$. For $q=5$: $\{0,4,6\}\equiv\{0,4,1\}$ misses $2$; $\{0,2,6\}\equiv\{0,2,1\}$ misses $3$; $\{0,2\}$ misses $1$.
\end{proof}

\section{Main results}\label{sec:main}

\begin{theorem}\label{thm:gaps}
For all $n\ge2$, $p_{n+1}-p_n\ge2$ and $p_{n+1}-p_n$ is even. Consequently $H_1\ge2$, and $H_1\le3$ if and only if $p_{n+1}-p_n=2$ for infinitely many $n$.
\end{theorem}
\begin{proof}
For $n\ge2$, $p_n$ and $p_{n+1}$ are odd and distinct, so their difference is a positive even integer. Hence a gap $\le3$ with $n\ge2$ equals $2$; the only gap equal to $1$ is $p_2-p_1$.
\end{proof}

\begin{theorem}\label{thm:main}
$\DHL[3,2]$ implies $H_1\le6$.
\end{theorem}
\begin{proof}
By \cref{lem:adm}, $\{0,4,6\}$ is admissible, so $\DHL[3,2]$ gives infinitely many $n$ such that at least two of $n,n+4,n+6$ are prime. For each such $n$ we obtain primes $p<q$ with $q-p\in\{2,4,6\}$ and $p\ge n$. Write $p=p_k$. Since $q$ is a prime larger than $p$, the next prime satisfies $p_{k+1}\le q$, so $p_{k+1}-p_k\le q-p\le6$. As $p\ge n$ and $n$ is unbounded, so is $k$.
\end{proof}

\begin{theorem}\label{thm:model}
Let $P(n)$ hold if and only if $n\bmod4\in\{0,1\}$. Then $\DHL_P[3,2]$ holds (indeed for every $3$-element set), but $P(n)\wedge P(n+2)$ holds for no $n$; in particular $\DHL_P[2,2]$ fails.
\end{theorem}
\begin{proof}
Integers differing by $2$ have residues differing by $2$ modulo $4$, so they are never both in $P$. If $d\not\equiv2\pmod4$, then $P(m)\wedge P(m+d)$ for all $m\equiv0\pmod 4$ when $d\equiv0,1$, and for all $m\equiv1\pmod4$ when $d\equiv3$. For $a<b<c$, the differences $b-a$ and $c-b$ cannot both be $\equiv2\pmod4$, since then $c-a\equiv0$. So some pair in the set has a difference $\not\equiv2\pmod 4$ and is hit infinitely often. Finally, $\{0,2\}$ is admissible by \cref{lem:adm}.
\end{proof}

\begin{remark}\label{rem:model}
\Cref{thm:model} shows only that no proof of $\DHL[3,2]\Rightarrow\DHL[2,2]$ can rely on the combinatorial form of $\DHL[3,2]$ alone. It is \emph{not} an independence result for the primes, and $\DHL[3,2]$ is compatible with $H_1=2$.
\end{remark}

\begin{hypothesis}[Parity barrier]\label{hyp:parity}
A given sieve method establishes $\DHL[k,j]$ only if $j\le\lceil k/2\rceil$.
\end{hypothesis}

This hypothesis reflects the parity discussion in~\cite{Polymath} for square-of-divisor-sum weights. It is assumed, not proved.

\begin{theorem}\label{thm:parity}
\leavevmode
\begin{enumerate}
\item Under \cref{hyp:parity}, a sieve method establishing $\DHL[k,2]$ has $k\ge3$.
\item (Unconditional.) No admissible $3$-element set is contained in an interval $[a,a+5]$. Together with \cref{lem:adm}, the minimal diameter of an admissible $3$-set is exactly $6$.
\end{enumerate}
\end{theorem}
\begin{proof}
(1) $2\le\lceil k/2\rceil$ forces $k\ge3$. (2) Admissibility modulo $2$ forces the three elements to share a parity. If they lie in an interval of length $5$, any two differ by $2$ or $4$, so they are pairwise distinct modulo $3$ and cover all residues modulo $3$, contradicting admissibility modulo $3$.
\end{proof}

\section{Formal verification}\label{sec:lean}

All results above are formalized in Lean~4 with Mathlib~\cite{Mathlib} in the repository \url{https://github.com/BenFrohman/DHL-3-2-implies-liminf-prime-gap-le-6}. The development contains no \texttt{sorry}; \texttt{\#print axioms} reports only \texttt{propext}, \texttt{Classical.choice} and \texttt{Quot.sound}. Primes are indexed from $0$ in Lean ($p_0=2$). Key definitions and statements (rendered here in ASCII notation; the source uses Lean's Unicode symbols):

\begin{lstlisting}
def H1Le (c : Nat) : Prop := forall N, exists n, n >= N /\ gap n <= c
def Admissible (H : Finset Nat) : Prop :=
  forall q, q.Prime -> exists r < q, forall h in H, h % q != r
def DHL (k j : Nat) : Prop :=
  forall H : Finset Nat, H.card = k -> Admissible H ->
    forall N, exists n, n >= N /\ j <= (H.filter (fun h => (n + h).Prime)).card

theorem dhl32_H1Le_six (h : DHL 3 2) : H1Le 6
theorem H1Le_three_iff_twin : H1Le 3 <-> forall N, exists n, n >= N /\ gap n = 2
theorem dhl32_not_implies_dhl22_combinatorially :
  exists (P : Nat -> Prop) (_ : DecidablePred P),
    DHLfor P 3 2 /\ (forall n, Not (P n /\ P (n + 2))) /\ Not DHLfor P 2 2
structure ParityBarrier (Proves : Nat -> Nat -> Prop) : Prop where
  bound : forall k j, Proves k j -> j <= (k + 1) / 2
\end{lstlisting}

{\raggedright
\noindent Correspondence with the text:
\begin{itemize}
\item \cref{lem:adm}: \texttt{admissible\_02}, \texttt{admissible\_026}, \texttt{admissible\_046}.
\item \cref{thm:gaps}: \texttt{two\_le\_gap}, \texttt{not\_H1Le\_one}, \texttt{H1Le\_three\_iff\_two}, \texttt{H1Le\_three\_iff\_twin}.
\item \cref{thm:main}: \texttt{dhl32\_H1Le\_six}.
\item \cref{thm:model}: \texttt{Ptoy\_DHL32}, \texttt{dhl32\_not\_implies\_dhl22\_combinatorially}.
\item \cref{thm:parity}: \texttt{ParityBarrier}, \texttt{parity\_two\_needs\_three}, \texttt{admissible\_three\_diam\_ge\_six}, \texttt{min\_diam\_admissible\_three}.
\end{itemize}\par}
Statements from the literature (the bound $246$, the conditional status of $\DHL[3,2]$, and the unverified Stadlmann claim) are not formalized.

\section{Discussion}\label{sec:discussion}

A proof of $H_1=2$ must go beyond the combinatorial content of $\DHL[3,2]$ (\cref{rem:model}) and, if it is sieve-theoretic, must overcome the limitation recorded in \cref{hyp:parity}. Arguments that overcome or sidestep parity-type limitations exist in other settings. Friedlander and Iwaniec~\cite{FI} showed that $X^2+Y^4$ represents infinitely many primes, using bilinear (Type~II) estimates to overcome the parity obstruction. Maynard~\cite{Maynard19} proved that there are infinitely many primes with a restricted digit, using Fourier-analytic control of Type~I and Type~II sums; this is a related use of bilinear-sum methods rather than a direct parity-breaking result of the Friedlander--Iwaniec kind. Until such an ingredient is brought to bear on prime gaps, $\DHL[3,2]$ and $H_1\le6$ remain the natural limit of the GPY/Maynard--Tao method under \cref{hyp:parity}.

\begin{thebibliography}{9}
\bibitem{FI} J.~Friedlander and H.~Iwaniec, \emph{The polynomial $X^2+Y^4$ captures its primes}, Ann.\ of Math.\ (2) \textbf{148} (1998), no.~3, 945--1040.
\bibitem{GPY} D.~A.~Goldston, J.~Pintz and C.~Y.~Y{\i}ld{\i}r{\i}m, \emph{Primes in tuples I}, Ann.\ of Math.\ (2) \textbf{170} (2009), no.~2, 819--862.
\bibitem{Maynard15} J.~Maynard, \emph{Small gaps between primes}, Ann.\ of Math.\ (2) \textbf{181} (2015), no.~1, 383--413.
\bibitem{Maynard19} J.~Maynard, \emph{Primes with restricted digits}, Invent.\ Math.\ \textbf{217} (2019), no.~1, 127--218.
\bibitem{Mathlib} The mathlib Community, \emph{The Lean mathematical library}, in: Proceedings of the 9th ACM SIGPLAN International Conference on Certified Programs and Proofs (CPP 2020), ACM, 2020, 367--381.
\bibitem{Polymath} D.~H.~J.~Polymath, \emph{Variants of the Selberg sieve, and bounded intervals containing many primes}, Res.\ Math.\ Sci.\ \textbf{1} (2014), Art.~12, 83~pp.
\end{thebibliography}

\end{document}
